
Three research communities have independently developed metrics for evaluating large language model (LLM) reliability: factuality evaluation measures whether models generate truthful content, adversarial robustness quantifies resistance to input perturbations, and calibration assesses confidence-accuracy alignment. Despite measuring seemingly related properties---a model's ability to ''know what it knows''---no unified framework connects these evaluation paradigms.

\subsubsection{The Problem}

LLMs fail in two primary modes relevant to deployment safety. First, they generate confident but incorrect information---so-called hallucinations---that users may trust due to fluent presentation. Second, they succumb to adversarial perturbations, producing different outputs for semantically equivalent inputs. Both failure modes undermine user trust and limit safe deployment in high-stakes domains.

These failures share a common characteristic: overconfident predictions on inputs where the model should express uncertainty. A model that generates hallucinations does so confidently rather than hedging. A model vulnerable to adversarial attacks fails to recognize that perturbed inputs lie outside its training distribution. This suggests that both failures may stem from poor calibration---the model's confidence does not reflect its actual accuracy.

\subsubsection{The Calibration Hypothesis}

We hypothesize that well-calibrated models, whose confidence scores align with their accuracy, can better detect both their own errors (enabling factuality) and unusual inputs (enabling robustness). If true, calibration quality (measured by Expected Calibration Error, ECE) should mediate the correlation between factuality metrics (TruthfulQA MC1) and robustness metrics (1 - Attack Success Rate).

Formally: Under the scope of open-weight LLMs evaluated on word-level adversarial perturbations, if a model exhibits higher factuality error detection accuracy (TruthfulQA MC1), then it will demonstrate higher adversarial robustness (1 - ASR on TextFooler), because calibration quality (lower ECE) provides a shared internal signal that enables both error detection and robustness.

\subsubsection{Gap in Existing Work}

Prior work has studied each dimension in isolation. TruthfulQA established benchmarks for factuality; TextFooler developed standard attacks for robustness; Minderer et al. analyzed calibration across architectures. However, no empirical study has quantified whether models that excel at factuality also resist adversarial attacks, or whether calibration explains this relationship.

\subsubsection{Contributions}

We propose a methodology to test the calibration-mediation hypothesis:

\begin{enumerate}
\item \textbf{Correlation Analysis}: Evaluate 12 open-weight LLMs on TruthfulQA MC1 (factuality) and TextFooler (robustness), computing Pearson correlation with bootstrap confidence intervals.
\end{enumerate}

\begin{enumerate}
\item \textbf{Confound Control}: Partial correlation controlling for model scale (log parameters), within-family analysis to control for training data effects.
\end{enumerate}

\begin{enumerate}
\item \textbf{Mediation Testing}: Baron-Kenny mediation analysis with ECE as the proposed mediator, testing whether calibration accounts for the factuality-robustness relationship.
\end{enumerate}

\begin{enumerate}
\item \textbf{Intervention Validation}: Temperature scaling experiments to test whether improving calibration improves both factuality and robustness.
\end{enumerate}

This paper presents our validated methodology and proof-of-concept results. Full experimental results remain incomplete due to simulated ASR values and limited model coverage.

